\section{Czego te twierdzenia nie mówią}\label{sec:czego-nie-mowia}
\subsection*{Nie każda teoria musi być niezupełna}\label{sec:19:1}
Istotny jest cały zestaw założeń: efektywność, niesprzeczność oraz odpowiednia siła arytmetyczna. Usunięcie choć jednego z nich może zmienić sytuację. Zbiór $\operatorname{Th}(\mathbb N)$ jest niesprzeczny i zupełny, ale nie jest rekurencyjnie przeliczalny. Teoria sprzeczna w logice klasycznej dowodzi każdego zdania, jest więc zupełna w technicznym sensie --- tyle że nie odróżnia żadnych odpowiedzi.

Inny prosty przykład stanowią wszystkie zdania prawdziwe w ustalonej skończonej strukturze o skończonym języku. Można algorytmicznie obliczać ich wartość: kwantyfikatory sprawdzamy po skończenie wielu elementach, a formuły według ich budowy. Ta teoria jest niesprzeczna, zupełna i rozstrzygalna. Nie zawiera jednak $Q$, której aksjomaty wymuszają nieskończenie wiele różnych numerałów. Nie ma więc konfliktu z twierdzeniem Rossera.

\subsection*{Nie wynika stąd, że dowody lub komputery są bezużyteczne}\label{sec:19:2}
Twierdzenia nie unieważniają żadnego poprawnego dowodu. Pokazują ograniczenia zakresu jednej ustalonej metody formalnej, nie bezwartościowość tego, co metoda rzeczywiście potrafi wykazać. Nie dowodzą, że każde interesujące pytanie jest niezależne, ani że każda próba automatyzacji musi zawieść.

Skończony dowód można sprawdzić, a wiele klas problemów można rozstrzygać algorytmicznie. Weryfikator formalny nie musi rozstrzygać wszystkich zdań, aby być użyteczny: jego zadaniem jest kontrola przedłożonego świadectwa w konkretnym rachunku. Twierdzenia nie obiecują też krótkich dowodów ani szybkiego przeszukiwania. Rozstrzygalność, istnienie dowodu i praktyczny koszt obliczeń to różne kwestie.

Nie wynika z nich również, że człowiek jest z konieczności nieomylny albo silniejszy od każdego możliwego algorytmu. Argument o prawdziwości zdania Gödla korzysta z założenia niesprzeczności konkretnego systemu; nie daje człowiekowi automatycznej wiedzy o niesprzeczności wszystkich systemów. Wnioski o świadomości, fizyce czy dowolnym potocznym „systemie przekonań” wymagałyby dodatkowych, osobno uzasadnionych założeń.

